\clearpage
\section{量測先行：driver 版本快的那 1\% 從哪裡來}
\label{sec:ab}

\begin{oneline}
第一冊量到 driver 版本比 UIO 版本快約 1\%，這次用四組實驗拆開，找到原因是 CPU 看 DMA buffer 的記憶體屬性從 Device 改成 Normal non-cacheable，讓每一幀「CPU 寫入輸入影像」那一步變快。
\end{oneline}

\begin{before}[原本怎麼做（第一冊第 10 節的結論）]
第一冊在同一顆 bitstream、同一次開機上比較，driver 版本平均快 0.8\%（另一次開機 1.0\%）。但 UIO 版本換成 driver 版本時，下面三件事是一起改的：
\begin{itemize}
\item \textbf{輪詢在哪裡做}：UIO 版本由 app 在 user space 讀 CTRL；driver 版本在 kernel 裡用 \code{readl_poll_timeout()}。
\item \textbf{buffer 在哪裡}：UIO 版本固定在 device tree 保留的 \code{0x6000_0000}；driver 版本從 CMA 配，這次在 \code{0x5010_0000}。
\item \textbf{CPU 端的記憶體屬性}：UIO 版本經 \code{/dev/mem} 映射（下面第 148--153 行），得到 Device-nGnRnE；driver 版本用 \code{dma_mmap_coherent()}，得到 Normal non-cacheable。
\end{itemize}
\snip{uio_devmem}
三件事一起變，所以第一冊只能寫「driver 版本沒有變慢」，不能說是哪一件讓它變快。
\end{before}

\begin{tech}
\begin{itemize}
\item \textbf{控制變因的 A/B 實驗}：四組設定（見下表）。第 1→2 組只換 buffer 在保留區裡的位置；第 2→3 組同時換了等硬體的地方（user space → kernel）與 buffer 的來源（保留區 → CMA）；第 3→4 組只換 CPU 端的記憶體屬性。每組跑 6 次，分兩輪、第二輪順序反過來，避免板子慢慢變熱或變慢時偏袒某一組。
\item \textbf{module parameter（模組參數）\code{map_device}}：載入 driver 時給 \code{map_device=1}，driver 改用 \code{pgprot_noncached()} 加 \code{remap_pfn_range()} 映射，得到和 \code{/dev/mem} 一樣的 Device 屬性。只給量測用。
\item \textbf{環境變數 \code{PL_ARENA_PA}}：讓 UIO 版本把 buffer 改放在保留區裡的另一個位置。
\item \textbf{\code{twn_bench} 第 3 部分}：直接量 CPU 寫入、讀取 buffer 的頻寬。
\item \textbf{\code{PIPE_TIMING=1}}：pipeline app 自己印出每一步花多少時間。
\end{itemize}
\end{tech}

\begin{gain}
\begin{itemize}
\item \textbf{找到原因}：只差記憶體屬性的兩組（第 3、4 組），fps 從 \abDrvDev{} 變成 \abDrvNc{}（+1.2\%）。前三組（\abUio{}、\abUioSix{}、\abDrvDev{}）相差不到 0.3 fps，所以輪詢位置與 buffer 位置合起來看不出影響；這兩件事在第 2→3 組是一起換的，沒有分開，也不能排除兩個效應正好互相抵銷。
\item \textbf{找到機制}：Normal non-cacheable 下，CPU 寫 buffer 的頻寬是 Device 的 3.1 倍（1-byte store）到 5.0 倍（8-byte store）；pipeline 每一幀「寫輸入」從 \splitDevIn{} ms 變成 \splitNcIn{} ms，差約 0.05 ms，和每幀週期的差（約 0.04 ms）相當（\splitDevPeriod{} 對 \splitNcPeriod{} ms）。
\item 第一冊留下的「分不出原因」可以結案：這 1\% 來自 driver 的映射方式，而且有機制上的解釋，不是量測誤差。
\end{itemize}
\end{gain}

\subsection{實際例子：四組實驗}

\begin{center}
\small
\begin{tabularx}{\textwidth}{l l l l X}
\toprule
組別 & 等硬體的地方 & buffer 位置 & CPU 端屬性 & 跟上一組只差 \\
\midrule
\code{uio} & app（user space） & \code{0x6000_0000} & Device-nGnRnE & （起點，v1.0.0 的做法） \\
\code{uio_68} & app（user space） & \code{0x6800_0000} & Device-nGnRnE & buffer 位置 \\
\code{drv_dev} & driver（kernel） & \code{0x5010_0000}（CMA） & Device-nGnRnE & 等硬體的地方與 buffer 來源（兩件一起換） \\
\code{drv_nc} & driver（kernel） & \code{0x5010_0000}（CMA） & Normal non-cacheable & 只有 CPU 端屬性 \\
\bottomrule
\end{tabularx}
\end{center}

每一組用 pipeline app 跑 120 幀（\code{FUSE_BB2HD=1}、head 在 PL），量 fps。實驗腳本的核心是 \code{run()}：先用第一冊的 \code{bind.sh} 把裝置交給對應的 driver，再跑三次：

\snip[bash]{ab_run}

\code{drv_dev} 這一組靠的是 driver 的 \code{map_device} 參數。\code{twn_mmap()} 在它設定時改走 \code{remap_pfn_range()}，屬性用 \code{pgprot_noncached()}，也就是 \code{/dev/mem} 給的 Device-nGnRnE；沒設定時走平常的 \code{dma_mmap_coherent()}：

\snip{drv_mmap}

\textbf{板上結果。}開機組 D（v1.0.0 的開機檔，v1.1 的 driver，這時 driver 是輪詢版）：

\boardfile[，第 27--31 行：每組 6 次的平均]{D1_ab_measure.log}{1-2,27-31}

\code{uio}、\code{uio_68}、\code{drv_dev} 三組都在 228.8 到 229.1 fps 之間，只有 \code{drv_nc} 跳到 \abDrvNc{} fps。各組 6 次的原始值在同一個檔案的第 3--26 行，組內的差距不超過 1 fps。

\textbf{為什麼記憶體屬性會影響 fps：直接量 CPU 存取 buffer 的速度。}\code{twn_bench} 第 3 部分在 buffer 的輸入區連續寫 1 MB、讀 1 MB：

\snip{bench_copy}

\boardfile[：第 11 行是預設的 Normal non-cacheable，第 20 行是 \code{map_device=1}]{D2_bench.log}{11-12,20}

寫入差很多（1-byte store \bwDevOne{} 對 \bwNcOne{} MB/s，8-byte store \bwDevEight{} 對 \bwNcEight{} MB/s），讀取一樣（都是 \bwLoad{} MB/s）。

\textbf{差異落在 pipeline 的哪一步。}pipeline app 開 \code{PIPE_TIMING=1}，會印出 producer（負責「寫輸入、啟動 PL、等 PL」的那個執行緒）最後一幀每一步的時間：

\boardfile[：前三次是 Normal non-cacheable，後三次是 \code{map_device=1}]{D6_stage_times.log}{3-21}

\code{write-in}（CPU 把這一幀的輸入影像寫進 buffer）從 \splitDevIn{} ms 變成 \splitNcIn{} ms，其他步驟（\code{backbone PL}、\code{head PL}）幾乎不變。每幀週期差約 0.04 ms，和 \code{write-in} 的差相當。寫輸入在 producer 的關鍵路徑上：寫完才能啟動 PL，所以它省下的時間直接變成 fps。

\subsection{原理}

\textbf{ARM 的兩種「不經過 cache」的記憶體。}arm64 的頁表每一頁都帶一個記憶體屬性。這次牽涉兩種，CPU 存取時都不經過 cache，差別在寫入能不能合併：

\begin{center}
\small
\begin{tabularx}{\textwidth}{l X X}
\toprule
 & Device-nGnRnE & Normal non-cacheable \\
\midrule
誰會得到 & \code{/dev/mem} 映射 no-map 區（第一冊第 3 節）；\code{pgprot_noncached()} & \code{dma_mmap_coherent()}（第一冊第 4 節）；\code{pgprot_dmacoherent()} \\
寫入 & 每一個 store 都是一筆獨立的匯流排交易，不合併（nG = non-Gathering）；寫入的回應必須由終點發出，中途的緩衝不能提早回應（nE = no Early write acknowledgement） & 連續的小 store 可以在 CPU 的 write buffer（CPU 送出寫入前暫放的小緩衝區）裡合併成一筆大的寫入 \\
讀取 & 每次都到 DDR & 每次都到 DDR（沒有 cache） \\
對齊 & 不能非對齊存取 & 可以 \\
\bottomrule
\end{tabularx}
\end{center}

所以寫入差 3 到 5 倍、讀取一樣，跟量到的完全對得上。讀取 81 MB/s 表示每次 8 bytes 的讀取大約要 100 ns，大約是一趟 DDR 讀取的延遲，兩種屬性都要付。

\textbf{為什麼只有 \code{write-in} 變快。}pipeline 的 CPU 對 buffer 做兩件事：寫輸入（store）、讀 heatmap（load）。store 受屬性影響、在關鍵路徑上；load 不受屬性影響，而且在另一個執行緒（consumer），和 PL 同時進行。所以只有寫輸入那一步的改善會反映到 fps。

\textbf{為什麼不乾脆用會 cache 的記憶體。}PL 經 HP port 存取 DDR，不會去看 CPU 的 cache（不是 coherent 的連接）。若 buffer 映射成 cacheable，CPU 寫完輸入後必須先把 cache 寫回 DDR，PL 寫完輸出後又必須先作廢 cache 再讀，driver 要提供 \code{dma_sync_single_for_device()} 這類同步介面，app 也要配合。這是可以做的下一步，這次沒有做。

\textbf{實驗設計的三個細節。}第一，兩輪的順序相反（\code{uio}→\code{drv_nc}，再 \code{drv_nc}→\code{uio}），板子若隨時間慢慢變快或變慢，兩輪的影響會互相抵銷。第二，\code{uio_68} 這一組是為了看 buffer 在 DDR 的位置會不會影響：同樣用 UIO，只換位置，fps 沒有變。第三，\code{uio_68}→\code{drv_dev} 一次換了兩件事（等硬體的地方、buffer 從保留區改成 CMA），要把它們也分開，還需要一組「driver 版、buffer 放在保留區」或「UIO 版、buffer 在 CMA」，這次沒有做；因為這三組的 fps 幾乎一樣，結論不受影響，但嚴格說只能講「兩者合起來看不出影響」。

\subsection{面試會問}

\begin{qa}{你說 driver 版本比較快，你怎麼知道是哪個改動造成的？}
driver 版本同時改了三件事：輪詢從 user space 搬到 kernel、buffer 從保留區改成 CMA、CPU 端的記憶體屬性從 Device 改成 Normal non-cacheable。我做了四組實驗，每組 6 次、兩輪順序相反：先只換 buffer 在保留區裡的位置，再換成 driver（輪詢位置和 buffer 來源一起換），最後只換記憶體屬性。前三組都在 229 fps 附近，只有最後一步從 \abDrvDev{} 跳到 \abDrvNc{}，所以原因是記憶體屬性；輪詢位置和 buffer 來源是一起換的，我只能說它們合起來看不出影響。再往下量機制：Device 屬性下每個 store 都是獨立的匯流排交易，Normal non-cacheable 可以在 write buffer 合併，CPU 寫 buffer 的頻寬差 3 到 5 倍；pipeline 裡「寫輸入」那一步從 0.14--0.15 ms 變成 0.09--0.10 ms，和每幀週期的差相當。讀取兩邊一樣慢，因為都不經過 cache。
\end{qa}

\subsection{關鍵函式速查}

\begin{center}
\small
\begin{tabularx}{\textwidth}{l l X}
\toprule
名稱 & 位置 & 做什麼 \\
\midrule
\code{map_device} & \code{twn_fuse.c}:\lineof{drv_param_map} & module parameter：映射成 Device 屬性（只給量測用） \\
\code{twn_mmap()} & \code{twn_fuse.c}:\lineof{drv_mmap} & 預設 \code{dma_mmap_coherent()}；\code{map_device} 時 \code{remap_pfn_range()} \\
\code{run()} & \code{test/ab_measure.sh}:\lineof{ab_run} & 切換 driver、跑 pipeline 三次 \\
第 3 部分 & \code{test/twn_bench.c}:\lineof{bench_copy} & CPU 寫、讀 buffer 的頻寬 \\
\code{v11_check.sh split} & \code{test/v11_check.sh}:\lineof{chk_split} & 兩種映射下的 pipeline 分段時間 \\
\bottomrule
\end{tabularx}
\end{center}
